\documentclass[11pt]{amsart}
\usepackage[T1]{fontenc}
\usepackage[utf8]{inputenc}
\usepackage{lmodern,microtype}
\usepackage[margin=27mm]{geometry}
\usepackage{amsmath,amssymb,amsthm,mathtools}
\usepackage{xcolor}
\usepackage[colorlinks=true,linkcolor=blue!50!black,citecolor=blue!50!black,urlcolor=blue!50!black]{hyperref}
\usepackage{xurl}
\setlength{\emergencystretch}{2em}

\newcommand{\workbenchurl}{https://github.com/KokunoYumeto/schemes-and-moduli-workbench}
\newcommand{\recordurl}{https://github.com/KokunoYumeto/mathematics-workbenches}

\newtheorem{theorem}{Theorem}
\newtheorem{thm}{Theorem}[section]
\newtheorem{lemma}[thm]{Lemma}
\newtheorem{proposition}[thm]{Proposition}
\newtheorem{corollary}[thm]{Corollary}
\theoremstyle{remark}
\newtheorem{remark}[thm]{Remark}

\newcommand{\B}{\mathbb B}
\DeclareMathOperator{\Spec}{Spec}
\DeclareMathOperator{\Idl}{Idl}
\DeclareMathOperator{\supp}{supp}

\title[Spectrum, semimodules and projectives of $R_0$]{The zero-adjoined ring $R_0$ of Faul, Goswami, Joubert and Manuell: spectrum, semimodules and projectives}
\author{}
\date{23 November 2026}

\begin{document}

\begin{abstract}
Faul, Goswami, Joubert and Manuell introduced inverse semirings, among them the semiring $R_0$ obtained from a commutative ring $R$ by adjoining a new zero that is an additive identity and a multiplicative absorber, and showed that every module over an inverse semiring has an inverse additive monoid. For a nonzero commutative ring $R$ we determine the prime spectrum of $R_0$ ($\Spec R$ with one new point below everything, whose residue semifield is the Boolean semifield), its congruences, and its semimodules: they are equivalent to diagrams of $R$-modules over join-semilattices. We then classify the finitely generated projective $R_0$-semimodules: their semilattices of idempotents are the finite distributive lattices, and their fibres are sums of projective $R$-modules indexed by order ideals. The classification rests on the difference between the old zero of $R$ and the new one. Finally, $R\mapsto R_0$ is fully faithful, so that $\operatorname{Aut}R_0=\operatorname{Aut}R$, and the two zeros are singled out by symmetry: when the units of $R$ have no common nonzero fixed point in $R$, the common fixed points of the units in $R_0$ are exactly the old zero and the new one, and they form the Boolean semifield. The spectrum, the congruences and special cases of the last section's facts go back to the project's first papers on the construction, written in November 2025 by Gemini~2.5 Deep Think (Gemini app, Deep Think mode); the semimodule and projective results were produced by GPT-5-series Pro models in ChatGPT Pro mode, at reasoning effort ultra, in its later notes on the ``split zero''. Claude Opus 5.5 at effort max checked all of them and wrote this text; no human has checked them.
\end{abstract}

\maketitle

\section{The construction and the results}

\subsection{Background}
An inverse semiring is a semiring whose additive monoid is an inverse monoid. Faul, Goswami, Joubert and Manuell made the case for this class between rings and semirings \cite{FGJM}. One of their basic constructions adjoins a zero: for an inverse semiring $R$, the inverse semiring $R_0=R\sqcup\{0\}$ has the operations of $R$ on $R$, and the new zero is an additive identity and a multiplicative absorber \cite[Definition 3.6]{FGJM}; it is the free addition of a zero, left adjoint to forgetting the zero \cite[\S3.1]{FGJM}, and the same adjunction appears for semirings in Golan's book \cite{Golan}. Their examples include $\mathbb Z_0$ and $\mathbb R_0$ \cite[Example 3.7]{FGJM}, and they prove that a module over an inverse semiring has an inverse additive monoid \cite[Proposition 6.1]{FGJM}. The same construction for commutative rings, under the name ``split zero'', is the starting point of this project's earlier notes.

\subsection{Notation}
Let $R\ne0$ be a commutative ring. To keep the two zeros apart we write $\tau$ for the adjoined zero and $e=0_R$ for the zero of $R$; so $G(R)=R_0=R\sqcup\{\tau\}$, with $\tau+x=x$ and $\tau x=\tau$ for all $x$. Then $e\ne\tau$, $e+e=e$ and $e^2=e$. Let $\B=\{0,1\}$ be the Boolean semifield, $1+1=1$.

\subsection{The theorems}

\begin{theorem}[spectrum and congruences]\label{thm:spec}
The ideals of $G(R)$ are $\{\tau\}$ and $J\cup\{\tau\}$ for the ideals $J$ of $R$, and its prime ideals are $\eta=\{\tau\}$ and $P\cup\{\tau\}$ for $P\in\Spec R$. Every point lies in the closure of $\eta$, $\{\eta\}$ is open, the closed set $V(e)$ is homeomorphic to $\Spec R$, and the localization at $\eta$ is $\B$; if $\dim R<\infty$, then $\dim G(R)=\dim R+1$. The congruences of $G(R)$ are, for each ideal $J$ of $R$, the one identifying $\tau$ with the class of $e$, with quotient $R/J$, and the one keeping $\tau$ alone, with quotient $G(R/J)$.
\end{theorem}

A \emph{support diagram} is a join-semilattice $L$ with least element, an $R$-module $A_l$ for each $l\in L$, and $R$-linear maps $\rho_{l,l'}\colon A_l\to A_{l'}$ for $l\le l'$ forming a functor.

\begin{theorem}[semimodules]\label{thm:mod}
For a $G(R)$-semimodule $M$ put $\varepsilon(m)=e\,m$, $L_M=eM$ and $M_l=\varepsilon^{-1}(l)$. Then $L_M$ is a join-semilattice under $+$, each $M_l$ is an $R$-module with zero $l$, the maps $T_{l,l'}(m)=m+l'$ form a support diagram, and $m+n=T(m)+T(n)$ in the fibre over $\varepsilon(m)+\varepsilon(n)$. This is an equivalence between $G(R)$-semimodules and support diagrams; the cancellative semimodules correspond to $R$-modules.
\end{theorem}

For a finite poset $J$ let $\Idl(J)$ be the distributive lattice of its order ideals, and for finitely generated projective $R$-modules $Q_j$ ($j\in J$) let $P(J,Q)$ be the support diagram over $\Idl(J)$ with fibre $\bigoplus_{j\in U}Q_j$ over $U$ and the summand inclusions as transports.

\begin{theorem}[projectives]\label{thm:proj}
A finitely generated $G(R)$-semimodule is projective if and only if it is isomorphic to some $P(J,Q)$. In particular the semilattices of idempotents of finitely generated projective $G(R)$-semimodules are exactly the finite distributive lattices.
\end{theorem}

\subsection{What the theorems say, and why they matter}
Theorem~\ref{thm:spec} says that $\Spec R_0$ is $\Spec R$ with one point added below everything, a generic point at which every element of $R$, including the old zero $e$, becomes the unit: the residue semifield there is $\B$. So adjoining a zero adds a ``characteristic one'' point to every affine scheme, in the sense in which $\B$ is the field of characteristic one of Connes and Consani \cite{CC}. Theorem~\ref{thm:mod} makes the module theory explicit: the inverse monoid of \cite[Proposition 6.1]{FGJM} is, by Clifford's theorem, a semilattice of groups, and over $R_0$ these groups are $R$-modules glued along a semilattice. Theorem~\ref{thm:proj} is the first classification we know of projective modules over a class of inverse semirings that are not rings: projectives are not just projective $R$-modules, but carry a finite distributive lattice, and every such lattice occurs. The difference between the two zeros is what produces it: the retraction that makes $P(J,R)$ projective fills order ideals with copies of $e$.

\subsection{Credit}
The construction is in \cite[Definition 3.6]{FGJM}, and the fact that modules over inverse semirings have inverse additive monoids is \cite[Proposition 6.1]{FGJM}; Clifford's description of commutative inverse monoids as strong semilattices of abelian groups is classical \cite{Clifford,Howie}; Birkhoff's representation of finite distributive lattices is used in Section~\ref{sec:proj} \cite{Birkhoff}. The spectrum, the congruences and special cases of Section~\ref{sec:sym} (the fixed points of the units for $\mathbb Z$ and rings of integers, zerosumfreeness and subtractive ideals for $\mathbb Z$, and $\operatorname{Aut}G(R)=\operatorname{Aut}R$) were found in this project's first papers on the construction, and the equivalence with support diagrams and the classification of projectives in its later notes on the split zero (see the note on method); all are proved here, those of Section~\ref{sec:sym} for every commutative ring and with full faithfulness in place of the automorphism group.

\subsection{Context}
Projective objects in the category of semilattices were determined by Horn and Kimura \cite{HK}, and projective distributive lattices by Balbes and Horn \cite{BH}; Lemma~\ref{lem:retract} below is consistent with their results in the finite case, and we prove it directly. Semirings with a negation map and their quasi-zeros $a+(-)a$ are studied by Rowen and coauthors \cite{Rowen}; in $G(R)$ the quasi-zero of $m$ is $m+(-1)m=\varepsilon(m)$, and Theorem~\ref{thm:mod} describes the modules fibred over their quasi-zeros. Lorscheid's blueprints and Connes and Consani's semirings of characteristic one are other frameworks for adjoining semiring points to schemes; we do not compare the constructions here.

\subsection{Scope}
Proved: Theorems~\ref{thm:spec}--\ref{thm:proj} for every nonzero commutative ring $R$, and the results of Section~\ref{sec:sym} for every commutative ring. Not treated: infinitely generated projectives; noncommutative $R$; the injective semimodules; the relation with Rowen's systems beyond the remark above.

\subsection{Proof overview}
Everything rests on two identities in a $G(R)$-semimodule: $m+e\,m=(1+e)m=m$ and $m+(-1)m=e\,m$, since $1+e=1$ and $1+(-1)=e$ in $R$. So $e\,m$ is an idempotent attached to $m$ that behaves like a local zero, and the semimodule falls apart into groups over the semilattice of these idempotents. For projectives, a retraction of a free module $G(R)^n$ restricts to a retraction of its semilattice of supports, the Boolean lattice of subsets of $[n]$, and the retracts of a finite Boolean lattice are exactly the finite distributive lattices; the fibres are retracts of free $R$-modules, and the lattice of fibres is distributive in the module sense, which splits them into the $Q_j$. Section~\ref{sec:sym} uses only that $G(R)$ is zerosumfree and has no zero divisors: a homomorphism cannot send the old zero to the new one, and a unit $u$ fixes an element of $R$ exactly when $u-1$ kills it.

\section{The spectrum and the congruences}

\begin{proof}[Proof of Theorem~\ref{thm:spec}]
The map $x\mapsto(p(x),\chi(x))$ with $p(\tau)=0$, $p|_R=\mathrm{id}$, $\chi(\tau)=0$, $\chi|_R=1$ is an injective semiring map $G(R)\to R\times\B$ with image $\{(x,b):b=0\Rightarrow x=0\}$; the operations are coordinatewise on the image.

An ideal $I\ne\{\tau\}$ contains some $r\in R$, hence $e=0\cdot r$, and then $I\cap R$ is an ideal of $R$ and $I=(I\cap R)\cup\{\tau\}$. A product $xy$ equals $\tau$ only if $x=\tau$ or $y=\tau$, so $\{\tau\}$ is prime and $P\cup\{\tau\}$ is prime exactly when $P$ is. Every prime other than $\eta$ contains $e$, so $\{\eta\}$ is the complement of $V(e)$, and $V(e)$ is homeomorphic to $\Spec R$ through $p$. As $\eta$ is contained in every prime, every point lies in its closure. Localizing at $\eta$ inverts every element of $R$, including the idempotent $e$; then $e=1$, so $r=r\cdot e$ becomes $1$ for every $r\in R$, and the localization is $\{\tau,1\}\cong\B$. A chain of primes of length $d$ in $R$ gives one of length $d+1$ in $G(R)$ by adding $\eta$ below it, and every chain arises so.

Let $\sim$ be a congruence. Its restriction to $R$ is a congruence of the ring $R$ (adding $-y$ to $x\sim y$ gives $x-y\sim e$), so it is $x\sim y\iff x-y\in J$ for an ideal $J$. If $\tau\sim r$ for some $r\in R$, then $\tau=\tau\cdot e\sim r\cdot e=e$, so the class of $\tau$ is $J\cup\{\tau\}$ and the quotient is $R/J$. Otherwise $\tau$ is alone in its class and the quotient is $G(R/J)$. Both relations are congruences.
\end{proof}

For a Dedekind domain $R=\mathcal O_K$ the finite quotients of the first kind are counted by $\sum_J N(J)^{-s}=\zeta_K(s)$, and those of the second kind have $N(J)+1$ elements. For $R=\mathbb Z$ the two families together give $\sum_{n\ge1}n^{-s}+\sum_{n\ge1}(n+1)^{-s}=2\zeta(s)-1$.

\section{Semimodules}

A $G(R)$-semimodule is a commutative monoid $(M,+,0_M)$ with an action of $G(R)$, distributive in both variables, with $1m=m$ and $\tau m=0_M$.

\begin{proof}[Proof of Theorem~\ref{thm:mod}]
In $R$, $1+e=1$ and $1+(-1)=e$, so $m+\varepsilon(m)=(1+e)m=m$, $m+(-1)m=\varepsilon(m)$, and $\varepsilon(m)+\varepsilon(m)=(e+e)m=\varepsilon(m)$. Hence each $M_l$ is an abelian group with identity $l$ and inverse $m\mapsto(-1)m$, $M$ is the disjoint union of the $M_l$, and $\varepsilon$ is additive, so $L_M=eM$ is a join-semilattice with least element $0_M=\varepsilon(0_M)$. This is Clifford's description of the commutative inverse monoid $(M,+)$ \cite{Clifford}, \cite[Ch.~4]{Howie}, with structure maps $T_{l,l'}$. Since $r\varepsilon(m)=\varepsilon(rm)$ for $r\in R$, the ring $R$ acts on each fibre, linearly, and the $T_{l,l'}$ are $R$-linear; $m+n=T(m)+T(n)$ in the fibre over the join is the identity $m+n=(m+l)+(n+l)$ for $l=\varepsilon(m)+\varepsilon(n)$. Conversely, a support diagram defines the set $\bigsqcup_lA_l$ with $x\boxplus y=\rho(x)+\rho(y)$ in the fibre over the join, $rx$ computed in the fibre for $r\in R$, and $\tau x=0\in A_{\bot}$; it is a $G(R)$-semimodule, and the two constructions are mutually inverse and functorial. A Clifford monoid is cancellative exactly when it has one idempotent, that is, when $L_M$ is a point.
\end{proof}

\section{Projectives}\label{sec:proj}

\begin{lemma}\label{lem:retract}
A finite join-semilattice with least element is a retract of a finite Boolean lattice $\mathcal P([n])$, by maps preserving joins and the least element, if and only if it is a distributive lattice.
\end{lemma}

\begin{proof}
If $L$ is such a retract via $i_0,r_0$, then $\epsilon=i_0r_0$ is an idempotent join-endomorphism of $\mathcal P([n])$ and $L\cong\epsilon(\mathcal P([n]))$ as posets. The image is closed under unions, and its meet is $U\wedge V=\epsilon(U\cap V)$: this lies below $U$ and $V$, and above every image point below both. As $\epsilon$ preserves unions, $U\wedge(V\cup W)=\epsilon((U\cap V)\cup(U\cap W))=(U\wedge V)\cup(U\wedge W)$. Conversely, by Birkhoff's theorem a finite distributive lattice is $\Idl(J)$ for a finite poset $J$ \cite{Birkhoff}, and $\Idl(J)\subseteq\mathcal P(J)$ is a retract through $S\mapsto{\downarrow}S$, which preserves unions and fixes order ideals.
\end{proof}

\begin{proof}[Proof of Theorem~\ref{thm:proj}]
\emph{If.} In $F=G(R)^J$ the support of $x$ is $\supp x=\{j:x_j\ne\tau\}$; then $\supp(x+y)=\supp x\cup\supp y$, and $\supp(rx)=\supp x$ for $r\in R$. So $P(J,R)=\{x:\supp x\in\Idl(J)\}$ is a subsemimodule, and it is the semimodule of the support diagram $P(J,(R)_j)$. Define $\rho\colon F\to P(J,R)$ by keeping $x$ on $\supp x$, putting $e$ on ${\downarrow}\supp x\setminus\supp x$ and $\tau$ elsewhere. Coordinatewise one checks $\rho(x+y)=\rho(x)+\rho(y)$ (on ${\downarrow}(S\cup S')\setminus(S\cup S')$ both sides are $e$, since $e+e=e+\tau=e$) and $\rho(rx)=r\rho(x)$ (as $re=e$); and $\rho$ is the identity on $P(J,R)$. So $P(J,R)$ is a retract of a free semimodule, hence projective. For fibres $R^{n_j}$ the same map on $G(R)^{J'}$, $J'=\bigsqcup_j\{j\}\times[n_j]$ ordered through $J$, is a retraction onto $P(J,(R^{n_j}))$. Finally, choosing $Q'_j$ with $Q_j\oplus Q'_j\cong R^{n_j}$, the fibrewise projections are maps of support diagrams and exhibit $P(J,Q)$ as a retract of $P(J,(R^{n_j}))$.

\emph{Only if.} A finitely generated projective $M$ is a retract of some $F=G(R)^n$ via semimodule maps $i,r$. They commute with $\varepsilon$, so they are maps of support diagrams, and $L_M$ is a retract of $L_F=\mathcal P([n])$; by Lemma~\ref{lem:retract}, $L_M\cong\Idl(J)$ for a finite poset $J$, with meet $l\wedge l'=r_0(i_0(l)\cap i_0(l'))$. The fibre $F_S$ over $S\subseteq[n]$ is $R^S$. Each $M_l$ is an $R$-module retract of $F_{i_0(l)}=R^{i_0(l)}$, hence finitely generated projective. Each transport $T_{l,l'}$ is split injective, with left inverse $r\circ\pi\circ i$ for the coordinate projection $\pi\colon R^{i_0(l')}\to R^{i_0(l)}$.

Write $A=i_0(l)$, $B=i_0(l')$. For $m\in M_{l\vee l'}$, split $i(m)\in R^{A\cup B}$ as $x+y$ with $x\in F_A$ and $y\in F_B$ (take $y$ to be $e$ on $A\cap B$); applying $r$ gives $M_{l\vee l'}=T(M_l)+T(M_{l'})$. If $z\in T(M_l)\cap T(M_{l'})$, then $i(z)$ equals $e$ outside $A\cap B$; let $w\in F_{A\cap B}$ agree with $i(z)$ on $A\cap B$ and be $\tau$ elsewhere. Then $w+\varepsilon_{A\cup B}=i(z)$, and since $i$ maps the idempotent $l\vee l'$ to $\varepsilon_{A\cup B}$, applying $r$ gives $z=T(r(w))$ with $r(w)\in M_{l\wedge l'}$. So $T(M_l)\cap T(M_{l'})=T(M_{l\wedge l'})$, and $l\mapsto T(M_l)$ maps the lattice $\Idl(J)$ to a lattice of direct summands of the top fibre, preserving joins and meets.

For $j\in J$ put $Q_j=M_{{\downarrow}j}/T(M_{{\downarrow}j\setminus\{j\}})$, a finitely generated projective module because the image of a split injection is a direct summand, and choose a section $s_j\colon Q_j\to M_{{\downarrow}j}$. Let $U\in\Idl(J)$ and $j_1,\dots,j_k$ a linear extension of $U$, so that $U_t=\{j_1,\dots,j_t\}$ are order ideals with $U_{t-1}\cap{\downarrow}j_t={\downarrow}j_t\setminus\{j_t\}$. By the two identities, $M_{U_t}=T(M_{U_{t-1}})+T(M_{{\downarrow}j_t})$ and $T(M_{U_{t-1}})\cap T(M_{{\downarrow}j_t})=T(M_{{\downarrow}j_t\setminus\{j_t\}})$, so $M_{U_t}/T(M_{U_{t-1}})\cong Q_{j_t}$. By induction, $(q_j)\mapsto\sum_{j\in U}T(s_j(q_j))$ is an isomorphism $\bigoplus_{j\in U}Q_j\to M_U$, compatible with the transports. So $M\cong P(J,Q)$ by Theorem~\ref{thm:mod}.

Every finite distributive lattice occurs, as the semilattice of idempotents of $P(J,R)$ for $J$ its poset of join-irreducibles.
\end{proof}

\begin{remark}
Over the ring $R$ alone, the skeleton disappears: the projective $R_0$-semimodules with one idempotent are the projective $R$-modules. The distributive lattice in Theorem~\ref{thm:proj} is created by the two zeros, since $\rho$ fills order ideals with the old zero $e$ and leaves $\tau$ outside them.
\end{remark}

\section{Functoriality and the symmetry of the new zero}\label{sec:sym}
The project's first papers on the construction proved that the two zeros are the fixed points of the units (for $\mathbb Z$ and for rings of integers), that $G(\mathbb Z)$ is zerosumfree with subtractive ideals, and that $\operatorname{Aut}G(R)=\operatorname{Aut}R$ for every commutative ring $R$. Here the first two facts are proved for every commutative ring, and the third in the stronger form of Proposition~\ref{prop:ff}. In this section the zero ring is allowed.

\begin{proposition}\label{prop:golan}
$G(R)$ is zerosumfree: $x+y=\tau$ only if $x=y=\tau$. It has no zero divisors: $xy=\tau$ only if $x=\tau$ or $y=\tau$. Every ideal $I$ of $G(R)$ is subtractive: if $a\in I$ and $a+b\in I$, then $b\in I$.
\end{proposition}

\begin{proof}
Sums and products of elements of $R$ lie in $R$, and $\tau$ is the only element outside $R$; this gives the first two statements. Let $a\in I$ and $a+b\in I$. If $b=\tau$ there is nothing to show, and if $a=\tau$, then $b=a+b\in I$. Otherwise $a$ and $a+b$ lie in $I\cap R$, which is closed under addition and under multiplication by elements of $R$; so $b=(a+b)+(-1)a$ lies in $I$.
\end{proof}

\begin{proposition}[full faithfulness]\label{prop:ff}
Let $R$ and $S$ be commutative rings. For a ring homomorphism $g\colon R\to S$ let $G(g)$ be $g$ on $R$ and send $\tau$ to $\tau$. Then $g\mapsto G(g)$ is a bijection from the ring homomorphisms $R\to S$ to the semiring homomorphisms $G(R)\to G(S)$ (preserving $0$ and $1$). In particular $\operatorname{Aut}G(R)=\operatorname{Aut}R$.
\end{proposition}

\begin{proof}
$G(g)$ is a homomorphism, since the operations involving $\tau$ are the same on both sides, and $g\mapsto G(g)$ is injective. Let $f\colon G(R)\to G(S)$ be a homomorphism; then $f(\tau)=\tau$. The element $f(e)$ satisfies $f(e)+f(e)=f(e)$, and in $G(S)$ the only such elements are $e_S=0_S$ and $\tau$, since $x+x=x$ forces $x=0$ in $S$. If $f(e)=\tau$, then $f(x)+f(-x)=f(e)=\tau$ for every $x\in R$, so $f(x)=\tau$ by Proposition~\ref{prop:golan}; in particular $f(1)=\tau$, which is not the unit of $G(S)$. So $f(e)=e_S$. Suppose $f(x)=\tau$ for some $x\in R$. Then $f(-x)=\tau+f(-x)=f(x)+f(-x)=f(e)=e_S$, so $\tau=f(x)=f(-1)f(-x)=f(-1)\,e_S$; this forces $f(-1)=\tau$, and then $f(1)=f(-1)^2=\tau$, again impossible. Hence $f(R)\subseteq S$, the restriction $g$ of $f$ to $R$ is a ring homomorphism, and $f=G(g)$.
\end{proof}

\begin{proposition}[the symmetry of the new zero]\label{prop:fix}
The units of $R$ act on $G(R)$ by multiplication, each by an automorphism of the additive monoid. Let $I$ be the ideal of $R$ generated by the elements $u-1$ with $u\in R^\times$. The common fixed points of the units are the elements of $\operatorname{Ann}_R(I)\sqcup\{\tau\}$. They are exactly the two zeros $e$ and $\tau$ if and only if $\operatorname{Ann}_R(I)=0$; this holds, for instance, when $R$ is a domain with a unit $u\neq1$, so for every ring of integers $\mathcal O_K$ (take $u=-1$) and every field other than $\mathbb F_2$. The pair $\{e,\tau\}$ is closed under both operations and is isomorphic to $\B$, with $e\mapsto1$ and $\tau\mapsto0$.
\end{proposition}

\begin{proof}
By distributivity each unit $u$ acts additively, with inverse $u^{-1}$. It fixes $\tau$, and it fixes $x\in R$ exactly when $(u-1)x=0$; so $x$ is fixed by every unit exactly when $Ix=0$. In a domain, $(u-1)x=0$ with $u\neq1$ forces $x=0$, and in $\mathcal O_K$ we have $-1\neq1$ because $2\neq0$. Finally, $e+e=e$, $e+\tau=e$, $\tau+\tau=\tau$, $e^2=e$, $e\tau=\tau$ and $\tau^2=\tau$ are the operations of $\B$.
\end{proof}

\begin{remark}\label{rem:chains}
For the zero ring, $G(0)=\{e,\tau\}$ with $e=1$ is the Boolean semifield $\B$. The support map $G(R)\to\B$, which sends $R$ to $1$ and $\tau$ to $0$, is $G$ applied to $R\to0$, and by Proposition~\ref{prop:ff} it is the only homomorphism $G(R)\to\B$. The construction of \cite[Definition 3.6]{FGJM} applies to every inverse semiring, in particular to the chain $L_n=\{0<1<\dots<n-1\}$ with $\max$ as addition and $\min$ as multiplication, whose additive monoid is a semilattice. Its new zero is a new least element, which absorbs for $\min$, so $(L_n)_0\cong L_{n+1}$: starting from $\B=G(0)\cong L_2$, adjoining zeros gives all the finite chains.
\end{remark}

\section{What remains open}
The classification of all projective $R_0$-semimodules without finiteness, the injective $R_0$-semimodules, and the extension of Theorems~\ref{thm:mod} and \ref{thm:proj} to the inverse semirings obtained from more general constructions of \cite[Construction 3.5]{FGJM} are not treated. Nor are the semimodules over the iterated extensions $(R_0)_0$, $((R_0)_0)_0$, \dots\ of Remark~\ref{rem:chains}. Whether the spectrum of Theorem~\ref{thm:spec}, glued over $\Spec\mathbb Z$, gives a useful ``characteristic one'' point in the sense of \cite{CC} is a question we leave open.

\section*{Note on method}
The construction for $\mathbb Z$ and for rings of integers, with its ideals, spectrum and congruences and the formula $\dim G(R)=\dim R+1$, was first written in November 2025 by Gemini~2.5 Deep Think (Gemini app, Deep Think mode), which formalized an idea from the project's conversation with Gemini~2.5 Pro (Gemini app); those papers were deposited on Zenodo under the name HI-AI and later withdrawn. In March 2026 a GPT-5-series Pro model in ChatGPT Pro mode, at reasoning effort ultra, extracted their usable part into a note on the ``globalization semiring'' $G(\mathbb Z)$, and GPT-5-series Pro models in ChatGPT Pro mode, at reasoning effort ultra, produced the statements and first proofs of the semimodule and projective results (Theorems~\ref{thm:mod} and \ref{thm:proj}) in the project's later notes. None of this work knew of \cite{FGJM}, which came first. Claude Opus 5.5 at effort max checked every proof (with exact computations of Theorem~\ref{thm:spec} for $G(\mathbb Z/n)$, $n\le12$, of the gluing $G(\mathbb Z/6)\cong G(\mathbb Z/2)\times_\B G(\mathbb Z/3)$, of Theorem~\ref{thm:mod} on 80 random subsemimodules of $G(\mathbb F_p)^n$, and of Lemma~\ref{lem:retract} for $n\le4$), wrote out the proof of Theorem~\ref{thm:proj} and the classification of congruences, proved the results of Section~\ref{sec:sym} in the generality given here, and wrote this text. No human has checked this note. Our text is dedicated to the public domain under CC0~1.0. The workbench is \expandafter\url\expandafter{\workbenchurl}; the permanent records are at \expandafter\url\expandafter{\recordurl}.

\begin{thebibliography}{9}
\bibitem{BH} R.~Balbes and A.~Horn, \emph{Projective distributive lattices}, Pacific J.\ Math.\ \textbf{33} (1970), 273--279.
\bibitem{Birkhoff} G.~Birkhoff, \emph{Rings of sets}, Duke Math.\ J.\ \textbf{3} (1937), 443--454.
\bibitem{Clifford} A.~H.~Clifford, \emph{Semigroups admitting relative inverses}, Ann.\ of Math.\ \textbf{42} (1941), 1037--1049.
\bibitem{CC} A.~Connes and C.~Consani, \emph{Homological algebra in characteristic one}, Higher Structures \textbf{3} (2019), 155--247; arXiv:1703.02325.
\bibitem{FGJM} P.~F.~Faul, A.~Goswami, G.~Joubert and G.~Manuell, \emph{The case for inverse semirings}, arXiv:2411.12318 (v1, 19 November 2024): Definition 3.6, Example 3.7 and \S3.1 (pp.~5--6), Construction 3.5, Proposition 6.1.
\bibitem{Golan} J.~S.~Golan, \emph{Semirings and their Applications}, Kluwer, 1999.
\bibitem{HK} A.~Horn and N.~Kimura, \emph{The category of semilattices}, Algebra Universalis \textbf{1} (1971), 26--38.
\bibitem{Howie} J.~M.~Howie, \emph{Fundamentals of Semigroup Theory}, Oxford University Press, 1995, Ch.~4.
\bibitem{Rowen} L.~H.~Rowen, \emph{Algebras with a negation map}, arXiv:1602.00353.
\end{thebibliography}

\end{document}
